\chapter{Python RL Training Pipeline}\label{ch:python}
\thispagestyle{pagestyle}

Getting a reinforcement learning agent to drive a virtual car is, in
principle, straightforward. The tricky part is wiring everything together
so the agent can actually talk to the simulation, watch what is happening
in real time, and resume a training run without losing work. This chapter
walks through how SHILATE does all three.

\section{Gymnasium Environment Wrapper}
\label{sec:gym-wrapper}

The Python side of SHILATE speaks Gymnasium \cite{Towers2024} --- the
standard interface that practically every modern RL library expects. The
wrapper class \texttt{ShilateEnv} lives in \texttt{environment.py} and
translates between the clean Gym API and the somewhat messier world of
MQTT messages flying back and forth with Unity.

The observation and action spaces are declared once in
\texttt{\_\_init\_\_} and never change at runtime:

\begin{lstlisting}[language=Python,
  caption={Observation and action space declarations from environment.py.},
  label={code:spaces}]
obs_size = ray_count + 2          # rays + normalised speed + normalised steer
self.observation_space = spaces.Box(
    low=0.0, high=1.0, shape=(obs_size,), dtype=np.float32
)
self.action_space = spaces.Box(
    low=np.array([-1.0, 0.0, 0.0], dtype=np.float32),
    high=np.array([ 1.0, 1.0, 1.0], dtype=np.float32),
    dtype=np.float32,             # [steer, throttle, brake]
)
\end{lstlisting}

Keeping everything in $[0,1]$ means the first layer of the network always
receives well-scaled inputs -- no normalisation wrapper needed.

The \texttt{step()} method is where the action-observation loop actually
happens. It sends the action to Unity over MQTT, then simply waits:

\begin{lstlisting}[language=Python,
  caption={The core step loop from environment.py.},
  label={code:step}]
def step(self, action):
    steer, throttle, brake = float(action[0]), float(action[1]), float(action[2])
    self._ctrl.send_action(steer, throttle, brake)

    self._ctrl.obs_event.clear()
    if not self._ctrl.obs_event.wait(timeout=OBS_TIMEOUT):
        # Unity stopped responding -- truncate rather than crash
        _emit("SHILATE-HEALTH",
              f"obs_timeout phase=step step={self._step_count}")
        return self._build_observation(), -1.0, False, True, {"reason": "obs_timeout"}

    obs    = self._build_observation()
    reward = self._ctrl.training_reward
    done   = self._ctrl.training_done
    self._step_count += 1
    return obs, reward, done, False, {}
\end{lstlisting}

\texttt{threading.Event.wait()} releases the Python GIL while blocking,
so the MQTT callback thread can still receive messages. If Unity goes
silent for more than five seconds --- say because the Editor was paused
or crashed --- the episode is marked \texttt{truncated=True} and
Stable-Baselines3 bootstraps the value from the last state instead of
applying a terminal penalty.

The MQTT connection lives in \texttt{controller.py}. The
\texttt{VehicleController} class maintains a \texttt{paho-mqtt} client
and exposes helpers like \texttt{send\_action()} and \texttt{set\_gear()}.
All topics are rooted at the prefix \texttt{env0/} --- a deliberate
single-environment design that keeps the code simple while the system is
still being developed.

\section{PPO Agent Configuration}
\label{sec:ppo-config}

The PPO agent is assembled in \texttt{train.py}. The actor and critic
networks are defined in \texttt{model.py} --- keeping them apart lets the
value function and the policy learn different representations of the same
input, which often helps in practice:

\begin{lstlisting}[language=Python,
  caption={Policy network architecture from model.py.},
  label={code:netarch}]
def get_policy_kwargs(ray_count: int = DEFAULT_RAY_COUNT) -> dict:
    return dict(
        net_arch=dict(
            pi=[256, 128],   # actor:  obs -> 256 -> 128 -> 3 actions
            vf=[256, 128],   # critic: obs -> 256 -> 128 -> scalar value
        )
    )
\end{lstlisting}

Back in \texttt{train.py}, the model is assembled and a checkpoint
callback saves progress every 10\% of the training budget. A resume path
can be supplied to pick up where a previous run left off:

\begin{lstlisting}[language=Python,
  caption={PPO model creation and checkpoint setup from train.py.},
  label={code:pposetup}]
if args.resume_model:
    model = PPO.load(args.resume_model, env=vec_env,
                     learning_rate=args.learning_rate,
                     n_steps=args.n_steps, batch_size=args.batch_size)
    model.tensorboard_log = args.log_dir
else:
    model = PPO("MlpPolicy", vec_env,
                learning_rate=args.learning_rate,
                n_steps=args.n_steps, batch_size=args.batch_size,
                policy_kwargs=get_policy_kwargs(ray_count),
                verbose=1, tensorboard_log=args.log_dir)

checkpoint_cb = CheckpointCallback(
    save_freq=max(total_timesteps // 10, 1000),
    save_path=args.save_path, name_prefix="shilate_ppo",
)
\end{lstlisting}

Default hyperparameters are listed in \cref{table:ppo-params}.

\begin{table}[H]
\caption{Default PPO hyperparameters.}
\label{table:ppo-params}
\begin{tabular}{|p{5cm}|p{3cm}|p{4cm}|}
  \hline
  \textbf{Parameter} & \textbf{Default} & \textbf{CLI flag} \\
  \hline
  Total timesteps      & 100{,}000          & \texttt{--total-timesteps} \\
  \hline
  Learning rate        & $3 \times 10^{-4}$ & \texttt{--learning-rate} \\
  \hline
  Rollout steps $n$    & 2048               & \texttt{--n-steps} \\
  \hline
  Batch size           & 64                 & \texttt{--batch-size} \\
  \hline
  Clip $\varepsilon$   & 0.2                & SB3 default \\
  \hline
  Discount $\gamma$    & 0.99               & SB3 default \\
  \hline
  GAE $\lambda$        & 0.95               & SB3 default \\
  \hline
\end{tabular}
\end{table}

\section{Training Loop and Health Telemetry}
\label{sec:training-loop}

One of the more unusual aspects of SHILATE is the structured telemetry
that the Python process streams back to the Unity Editor. Instead of just
printing progress to a terminal, every meaningful event emits a tagged
line that the Editor's parser picks up and displays in the Training
Controller window in real time.

The \texttt{HealthCallback} class runs at the end of each rollout. It
re-emits Stable-Baselines3 internal metrics as \texttt{SHILATE-METRIC}
lines, and checks loss values for NaN or infinity --- a surprisingly
useful guard during early experimentation with learning rates:

\begin{lstlisting}[language=Python,
  caption={Health callback emitting metrics and checking for non-finite losses.},
  label={code:healthcb}]
def _dump_metrics(self) -> None:
    stats = getattr(self.model.logger, "name_to_value", {}) or {}
    for key in self._WATCH:
        if key in stats:
            val   = stats[key]
            short = key.split("/", 1)[-1]
            emit("SHILATE-METRIC", f"{short}={val}")
            if "loss" in key or "kl" in key:
                if math.isnan(float(val)) or math.isinf(float(val)):
                    emit("SHILATE-HEALTH", f"nan_loss key={key} value={val}")
\end{lstlisting}

The \texttt{SHILATE-HEALTH} lines feed into
\texttt{TrainingHealthEvaluator.cs} in the Editor, which updates the
health banner colour accordingly. Alongside periodic checkpoints, the
final model is saved on clean exit --- or on \texttt{KeyboardInterrupt},
because most training runs are stopped manually rather than running to
the full timestep limit.

\section{Inference Mode}
\label{sec:inference}

Running a trained policy is handled by \texttt{ai\_driver.py}, which
loads a \texttt{.zip} checkpoint and feeds live observations into it.
The key difference from training is \texttt{deterministic=True}: this
disables the Gaussian exploration noise added during training, so the
policy always picks the action with the highest predicted return. In
plain terms --- the agent stops experimenting and drives with what it has
learned.

The inference loop itself is intentionally simple. It mirrors the same
\texttt{obs\_event.wait()} pattern used in training, but without the
Gymnasium wrapper overhead, and it handles episode resets automatically:

\begin{lstlisting}[language=Python,
  caption={Inference loop from ai\_driver.py.},
  label={code:aidriver}]
while _running:
    obs = build_observation(ctrl, ray_count)

    action, _ = model.predict(obs, deterministic=True)
    steer, throttle, brake = float(action[0]), float(action[1]), float(action[2])
    ctrl.send_action(steer, throttle, brake)
    step += 1

    ctrl.obs_event.clear()
    if not ctrl.obs_event.wait(timeout=2.0):
        log.warning("Observation timeout at step %d", step)
        continue

    if step % 50 == 0:
        log.info("Step %d | SPD: %.1f km/h | STR: %+.2f | THR: %.2f | BRK: %.2f",
                 step, ctrl.get_speed(), steer, throttle, brake)

    if ctrl.training_done:
        episode += 1
        log.info("Episode %d complete (step %d)", episode, step)
        step = 0
        time.sleep(0.5)
        ctrl.send_reset()
        ctrl.obs_event.clear()
        ctrl.obs_event.wait(timeout=5.0)
\end{lstlisting}

Running the inference loop on the deployed Leda image requires no changes
to the code --- only the \texttt{--model} path argument changes. This is
the same container image used during development, rebuilt and pushed to
the device registry by \texttt{provision-device.sh}.

\section{Cross-Platform Process Management}
\label{sec:process-mgr}

Because Unity typically runs on Windows while the Python process runs
inside WSL2, \texttt{PythonProcessManager.cs} handles virtual-environment
activation differently per platform. On Windows it calls
\texttt{Scripts\textbackslash{}activate}; on Linux/macOS it uses
\texttt{bin/activate}. Graceful shutdown sends SIGTERM to the process
group so TensorBoard logs are flushed and the final checkpoint is written
before Play mode exits. A hard kill is used as a fallback if the process
does not respond in time.
